\section*{How to Prepare Mathematical Notes}
\addcontentsline{toc}{section}{How to Prepare Mathematical Notes}

A mathematical note does not merely record an answer. It shows how the problem
is understood, which information matters, why a method is appropriate, how each
important step is justified, and how the conclusion follows from definitions,
assumptions, and earlier results.

Mathematical expressions are compact. A single formula may contain objects,
operations, conditions, and logical relationships. Reading a formula therefore
means being able to explain what each object represents, why the expression is
meaningful, and what information it provides.

A sequence of formulas without explanatory sentences rarely communicates a
complete argument. Words connecting the formulas are part of the mathematics
because they state the logical relationship between successive steps.

\subsection*{Make the Note Self-Contained}

Prepare one separate note for each assigned problem, including all subparts.
Begin with the problem or a complete and faithful formulation of it. Record what
is given, what must be found or proved, and what notation will be used.

Before calculating, identify the mathematical objects involved and decide what
would count as a satisfactory answer. After several weeks, opening the note
should still reveal both the question and the full solution without searching
elsewhere.

\subsection*{Make the Reasoning Visible}

Sentences that reveal mathematical reasoning include:
\begin{itemize}
    \item Because the determinant is nonzero, the matrix is invertible.
    \item We may therefore multiply both sides by the inverse matrix.
    \item By the definition of the derivative, we consider the following limit.
    \item Since the direction vector is nonzero, it determines a line.
    \item Substituting into the original equation verifies the solution.
\end{itemize}

The amount of explanation should be proportional to the problem. A simple
exercise may need only a few sentences; a substantial argument may need many.
The goal is not maximal length but complete visibility of every essential step.

Write as briefly as possible, but as completely as necessary.

\subsection*{A Complete Solution}

\begin{enumerate}
    \item \textbf{Understand the problem.}
    Record the question, the given information, and the required conclusion.

    \item \textbf{Choose a method.}
    Select definitions, theorems, or techniques connecting the data with the
    goal.

    \item \textbf{Construct the argument.}
    Carry out the necessary steps and explain why each important transition is
    valid.

    \item \textbf{Verify the result.}
    Check calculations, assumptions, domain restrictions, or geometric meaning.

    \item \textbf{State the conclusion.}
    Answer the original question clearly and in context.
\end{enumerate}

A bare final answer, an unexplained sequence of formulas, or disconnected
calculations constitute incomplete work. Quality cannot be replaced by quantity,
but a single polished solution does not replace completing the rest of the
assigned problem set.

\subsection*{Working with AI}

AI may clarify notation, generate simpler examples, check reasoning, identify
missing steps, and help organise mathematical prose. The student remains
responsible for every definition, calculation, assumption, conclusion, and
formula included in the submitted note.

Receiving a polished answer is not the same as understanding it. A note is
insufficient if the student cannot read its formulas, explain its terminology,
justify its transitions, or present the argument independently.

Work on one problem at a time and prepare one complete note for it. Do not paste
an entire problem set into one conversation and request all answers at once.
Repeat the focused process separately for every problem in the assigned set.

Useful questions to ask while working with AI include:
\begin{itemize}
    \item What does this formula mean?
    \item What mathematical object does each symbol represent?
    \item Why is this step valid?
    \item Which definition or theorem is being used?
    \item Are the assumptions of the theorem satisfied?
    \item Is there a simpler explanation of this transition?
    \item How can the result be checked?
    \item What should I be able to explain during class?
\end{itemize}

Preparing notes in this way develops the ability to read notation, recognise
objects, organise information, construct arguments, and communicate ideas
precisely. Repeating the process for the complete problem set builds habits of
systematic academic and professional work.

\subsection*{Student GitHub Workspace}

The current instructions, problem templates, and technical workflow are
maintained in the public repository
\href{https://github.com/dchorazkiewicz/Math_Problems_Repo}
{\texttt{dchorazkiewicz/Math\_Problems\_Repo}}. Students should create a public
fork and commit notes to their own fork; the common course repository is a
template, not the submission location.

\begin{itemize}
    \item \texttt{README.md} explains the workspace and weekly index;
    \item \texttt{SOURCE\_MATERIAL.md} maps weeks to the current PDF and LaTeX
          problem sources;
    \item \texttt{NOTE\_GUIDELINES.md} defines a complete note;
    \item \texttt{AI\_WORKFLOW.md} describes focused work, review, commits,
          presentation, feedback, and revision;
    \item \texttt{AGENTS.md} gives connected assistants operational rules;
    \item \texttt{docs/} contains weekly folders and one Markdown file per
          problem.
\end{itemize}
The workspace uses Markdown, Git and GitHub, MathJax, MkDocs, and GitHub Actions.
Because the technical instructions may change independently of this PDF,
students should consult the repository for the current workflow.
